\section{Demonstration}
\label{sec:demo}

The booth demo is a 10-minute arc through one module, mirroring how the
course is used.

\paragraph{(1) The whole lifecycle in 30 seconds.} We run
\texttt{course/00-the-map/journey.py}: a tiny model is built, trained on a
toy corpus, and sampled---loss falls from $\sim$6.0 to $\sim$2.5 live. This
frames everything else: the rest of the course makes each leg of the loop
\emph{good}.

\paragraph{(2) Read a module.} On the rendered site we open Module 03 (MLA):
the wall (KV memory), the theory-to-code table (each term of
$c = xW_{DKV}$, cache $c$, $K,V = \mathrm{split}(cW_{UKV})$ mapped to the
exact lines in \texttt{attention.py}), and the pasted, drift-guarded source.

\paragraph{(3) Build it.} The visitor solves
\texttt{lab\_mla.py}---two matmuls---and watches the grader compare their
latent round-trip against the real module (Figure~\ref{fig:session}):
\texttt{PASS} on success, a targeted assertion message on failure. For
skeptical visitors we rename a function in \texttt{attention.py} and run the
drift guard to show the documentation fail loudly.

\begin{figure}[t]
\begin{lstlisting}
$ python course/03-attention-mla/lab_mla.py
NotImplementedError: implement me - delete
  this line and write the two matmuls
# edit: latent = kv @ w_down
#       reconstructed = latent @ w_up
$ python course/03-attention-mla/lab_mla.py
PASS -- you implemented the core of MLA.
$ python course/03-attention-mla/ablation.py
KV cache, bytes/token/layer (bf16):
  MHA (n_kv_head=8)   2048  1.0x smaller
  MQA (n_kv_head=1)    256  8.0x smaller
  MLA (latent=128)     256  8.0x smaller
\end{lstlisting}
\caption{\textbf{Session snapshot of demo steps (3)--(4)}, verbatim: the lab
fails until implemented, the grader confirms against the real module, and the
module's ablation prints its measured payoff.}
\label{fig:session}
\end{figure}

\paragraph{(4) Measure and break it.} We run the module's ablation (the
$8\times$ cache table), then halve \texttt{kv\_lora\_rank} in a preset and
re-run---the number moves, the tradeoff is felt. Finally we start the
continuous-batching server, submit two concurrent chat requests, and show the
short one finishing first while both stream. Because the target platform is
the machine many attendees carry, the demo doubles as an install-fest:
visitors can clone the repository and reproduce every step on their own Mac
during the session.

\paragraph{Availability.} Code, course, and CI are public under MIT:
\url{https://github.com/jc-su/BabyWhale} (rendered course:
\url{https://jc-su.github.io/BabyWhale}). Requirements: an Apple-Silicon Mac
(MLX) and Python 3.14; every demo step above also runs in the project's CI.
